\chapter{Consciousness: From Mathematical Admissibility to Phenomenal Occupancy}

\section{The Central Asymmetry}

The mathematical work developed so far demonstrates how states can remain equivalent under one form map while differing in richer descriptions, how representation theory can exclude symmetries that are impossible in a given sector, and how higher-order invariants can resolve distinctions unavailable at lower order. Consciousness introduces a new demand. It is not enough to construct a mathematical state with properties that resemble the structure of experience. The question is whether phenomenal experience is actually present in, identical with, generated by, or otherwise related to that structure.

This yields the third governing principle of the book:

\begin{principle}[Admissibility is not occupancy]
The existence of a mathematically coherent representation capable of carrying properties associated with experience does not establish that phenomenal experience occupies that representation.
\end{principle}

The principle is structurally analogous to the Platonic admissibility result. If an invariant multiplicity vanishes, the form is impossible in that representation. If the multiplicity is nonzero, the form is allowed, not actualized. Consciousness theories face the same asymmetry at a more difficult semantic level. A model whose state space cannot represent a property it claims to explain can be ruled out. A model whose state space can represent it has merely survived an exclusion test.

The mistake is especially tempting when the mathematical structure is beautiful. Integration, recursion, self-reference, symmetry, nonseparability, projective geometry, or high-dimensional relational complexity may all be relevant to consciousness. None of these becomes phenomenality by being given a suggestive name. Mathematics can constrain a bridge theory; it cannot manufacture the bridge's destination.

\section{Why Reports Are Not Experience}

Every empirical access route to consciousness produces physical data. Verbal reports are acoustic or textual events. Button presses are motor events. Confidence ratings are behaviors. Eye movements, pupil responses, autonomic signals, scalp potentials, field potentials, spikes, hemodynamic changes, and perturbational responses are physical measurements. No-report paradigms can reduce contamination from explicit report preparation, but they remain dependent on physical indicators and can introduce their own inferential problems. The absence of report does not transform a measurement into direct access to phenomenality.

We therefore reserve $E_{\mathrm{exp}}$ for the experiential state hypothesized by an ontology and introduce a different symbol for the empirical construct:

\[
Z_\varphi=\text{latent phenomenological-evidence structure}.
\]

The permanent inequality is conceptual rather than numerical:

\[
Z_\varphi\neq E_{\mathrm{exp}}
\]

unless an additional occupancy argument is supplied. Multiple observable channels may be modeled as

\[
R_i=h_i(Z_\varphi,M_i,\epsilon_i),
\]

where $M_i$ represents channel-specific influences and $\epsilon_i$ noise. The purpose of combining verbal judgments, forced choice, confidence, involuntary eye behavior, pupil response, autonomic state, and other indicators is not to make $Z_\varphi$ nonphysical. It is to test whether a common latent structure survives across causally different output pathways.

Cross-channel stability is stronger than a single-report correlation because it makes simple output-specific confounds less plausible. Yet even perfect cross-channel reconstruction would still produce an empirical model of phenomenological evidence rather than an unmediated observation of experience. The semantic ceiling remains.

\section{An Evidential Ladder}

The gap between mathematical possibility and ontology should not be crossed in one step. We therefore organize the research program into levels. At the first level, formal admissibility asks whether the proposed structure is mathematically coherent and capable of representing the claimed relations. Structural nonredundancy then asks whether it is genuinely new rather than a reparameterization of an existing variable. Predictive necessity asks whether it improves held-out prediction. Interventional necessity asks whether perturbations change outcomes as the theory predicted before the intervention. Cross-context invariance asks whether approximately the same relation survives new tasks, subjects, modalities, and states. Independent phenomenological anchoring asks whether the phenomenological evidence was constrained independently of the physical theory against which it will be tested. Only then does it become meaningful to compare empirical occupancy models, and even the best occupancy model does not automatically settle ontology.

We label the stages $A_0$ through $A_7$:

\[
A_0:\ \text{formal admissibility},
\]

\[
A_1:\ \text{structural nonredundancy},
\]

\[
A_2:\ \text{predictive necessity},
\]

\[
A_3:\ \text{interventional necessity},
\]

\[
A_4:\ \text{cross-context invariance},
\]

\[
A_5:\ \text{independent phenomenological anchoring},
\]

\[
A_6:\ \text{best-supported empirical occupancy model},
\]

\[
A_7:\ \text{ontological interpretation}.
\]

The relation among them is deliberately non-implicational:

\[
A_0\nRightarrow A_1\nRightarrow A_2\nRightarrow A_3\nRightarrow A_4\nRightarrow A_5\nRightarrow A_6\nRightarrow A_7.
\]

Each transition requires additional evidence. The ladder prevents a common pattern in speculative consciousness theory in which a mathematical construction begins at $A_0$ and its conceptual resemblance to experience is treated as though the argument had already reached $A_7$.

\section{The Hard Problem as Several Ceilings}

The phrase ``hard problem'' becomes more useful when decomposed. The representation ceiling asks what mathematical object, if any, adequately represents phenomenological structure. The identifiability ceiling asks whether such structure can be inferred from data. The causal ceiling asks whether it contributes anything beyond an increasingly comprehensive physical description. The semantic ceiling asks why the recovered structure should be interpreted as experience rather than another physical latent variable. The ontological ceiling asks whether even empirical indispensability would establish that interiority is irreducible.

The present mathematical program has the greatest leverage on the first three. Representation theory can exclude inadequate structures. Model comparison can test identifiability and predictive necessity. Interventions and dynamical closure can test causal structure. The semantic and ontological ceilings remain more difficult because they concern what the successful formal variable is, not merely what relations it predicts.

This division does not trivialize the hard problem. It locates it. Progress can occur even if $A_7$ remains unreachable. A theory can eliminate impossible bridge structures, improve phenomenological measurement, identify causal invariants, or show that a proposed ontology makes no empirical difference. The success condition of the program is therefore not limited to proving a final metaphysics. Better constraints and eliminated theories are genuine scientific gains.

\section{Related Consciousness Programs}

Existing physicalist consciousness theories can be understood as different attempts to make $P_{\mathrm{phys}}$ sufficient. Global workspace approaches emphasize large-scale availability and ignition; integrated information theory begins from axiomatized phenomenal properties and proposes corresponding intrinsic cause--effect structure; predictive-processing and active-inference approaches formulate consciousness within hierarchical generative dynamics; higher-order theories emphasize representation of mental states by other mental states. These programs differ substantially, but all can be evaluated using the same separation between admissibility, predictive necessity, and occupancy.

The distinction is especially sharp for theories that derive mathematical structures from phenomenological axioms. A theory may successfully derive a formal object from a set of stated phenomenal properties while leaving open whether biological systems actually instantiate the object, whether the object is necessary for experience, and whether the formal structure is identical with experience rather than correlated with it. This is precisely the admissibility--occupancy boundary.

No-report paradigms offer a complementary methodological lesson. They attempt to reduce confounds associated with explicit report, but they remain dependent on proxies whose validity must itself be established. They should therefore be understood as confound-ablation strategies inside the phenomenological-evidence program rather than as direct access to consciousness.

The next chapter turns these principles into a concrete empirical program designed to make physical sufficiency and relational complementarity lose or win under preregistered conditions rather than philosophical preference.
